\subsection{单表示的维数}

用 n 表示群 G 的基数。设 \(\pi\) 是 G 在有限维 K-向量空间 M 中的一个线性表示。对每个 \(g \in G\)，我们有 \(\pi(g)^{n} = 1_{M}\)，故 \(\pi(g)\) 的最小多项式整除 \(T^{n} - 1\)。由于在 K 中 \(n \cdot 1 \neq 0\)，这个最小多项式是可分的（V, §11, No.~2, 第 78 页）；又由于体 K 是代数闭的，M 的自同态 \(\pi(g)\) 可对角化（VII, §5, No.~7, 第 40 页，命题 12）。\(\pi(g)\) 的特征值都是 n 次单位根，并且对每个 \(\alpha \in K\)，\(\alpha\) 作为 \(\pi(g)\) 的特征值的几何重数（VII, §5, No.~2, 第 30 页，定义 1）等于 \(\alpha\) 作为 \(\pi(g)\) 的特征多项式的根的重数。用 \(\mathcal{O}_{n}\) 表示由 n 次单位根的集合 \(\mu_{n}(K)\) 生成的 K 的子群；它是一个有限生成的 Z-模，并且是 K 的一个子环。\(\pi\) 的特征标取值于 \(\mathcal{O}_{n}\)。

命题 9. —— 假设体 K 的特征为零。那么 G 的每个单表示的次数都整除 G 的基数。

设 \((\mathrm{V},\pi)\) 是 G 的一个单表示，\(\chi\) 是它的特征标。对 \(\mathrm{Z}(\mathrm{K}[\mathrm{G}])\) 的每个元素 a，V 的自同态 \(\pi(a)\) 是一个位似（VIII, 第 47 页，定理 1）；用 \(\varphi(a)\) 表示使得 \(\pi(a)=\varphi(a)_{\mathrm{V}}\) 的标量。

这样得到的映射 \(\varphi\)（从 \(Z(K[G])\) 到 \(K\) 的映射）是一个代数同态。取 \(a = \sum_{g \in G} \chi(g^{-1})g\)；由 VIII, 第 408 页的注，我们有 \(\varphi(a) = (\dim V)^{-1}|G|\)。另一方面，\(a\) 属于 \(\mathcal{O}_n[G] \cap Z(K[G])\)，即 \(K[G]\) 的一个子环，后者是一个有限生成的 \(\mathbf{Z}\)-模（VII, §3, 第 15 页，推论）。因此元素 \(\varphi(a) = (\dim V)^{-1}|G|\)（\(K\) 的元素）属于 \(K\) 的一个作为有限生成 \(\mathbf{Z}\)-模的子环。我们借助下面的引理来完成证明。

引理. —— 设 L 是 Q 的一个扩张，A 是 L 的一个子环。假设 A 是有限生成的 Z-模。我们有 \(A \cap Q = Z\)。

由于 Z-模 \(A \cap Q\) 是有限生成的，存在严格正整数 N，使得 \(A \cap Q\) 包含于 \(\frac{1}{N}Z\)。设 x 是 \(Q - Z\) 的一个元素；把它写成 \(x = \frac{p}{q}\)，其中 p 与 q 是互素的整数且 \(q \geqslant 2\)。我们有 \(q^{N} \geqslant 2^{N} > N\)（集合论 III, §3, No.~6, 第 165 页，定理 2），整数 \(p^{N}\) 与 \(q^{N}\) 互素，从而 \(x^{N} \notin \frac{1}{N}Z\)。由此可知 x 不属于 A。引理证毕。

在下一小节中，我们把命题 9 推广到只假设 K 的特征不整除 G 的阶的情形。

